% TOP_PROBLEM: {"id": 30006887, "title": "Kobayashi Conjecture: Hyperbolic Compact Kähler Manifolds Have Ample Canonical Bundle", "category_id": 5, "category": {"id": 5, "name": "algebraic_geometry", "display_name": "Algebraic Geometry", "description": "Geometric objects defined by polynomial equations.", "slug": "algebraic-geometry", "order_index": 5, "created_at": "2026-07-31T15:26:25.671Z"}, "status": "open"}
% FIRST_POSTED: 2026-09-25
% ENTRY_KIND: statement_only
% !TeX program = lualatex
% Definition and sources only; this file is not an LLM solution attempt.
\documentclass[11pt]{article}
\usepackage[margin=25mm]{geometry}
\usepackage{fontspec}
\setmainfont{Latin Modern Roman}
\usepackage{amsmath,amssymb,mathtools}
\usepackage{unicode-math}
\setmathfont{Latin Modern Math}
\usepackage{xurl,hyperref}
\hypersetup{colorlinks=true,urlcolor=blue}
\setcounter{secnumdepth}{0}
\setlength{\parindent}{0pt}
\setlength{\parskip}{5pt}
\setlength{\emergencystretch}{3em}
\begin{document}
\section*{\texorpdfstring{Kobayashi Conjecture: Hyperbolic Compact Kähler Manifolds Have Ample Canonical Bundle}{Kobayashi Conjecture: Hyperbolic Compact Kähler Manifolds Have Ample Canonical Bundle}}\label{30006887}
\subsection{Definitions and mathematical statement}
A compact connected complex manifold is K\"ahler if it has a Hermitian metric whose associated real $(1,1)$-form is closed. The Kobayashi pseudodistance is the largest pseudodistance contracted by all holomorphic maps from the unit disk with its Poincar\'e metric. Hyperbolicity means it separates distinct points; in the compact setting this is equivalent to every holomorphic map ${\mathbb{C}}\to X$ being constant. The conjecture asserts
\[
 X\text{ compact K\"ahler and Kobayashi hyperbolic}
 \quad\Longrightarrow\quad K_X\text{ ample}.
\]
Here $K_X=\bigwedge^{\dim_{{\mathbb{C}}} X}\Omega_X^1$, and ample means that a positive power gives a projective embedding. Projectivity is thus part of the proposed conclusion, not an extra hypothesis. Negative curvature of a chosen metric is not assumed.
\par\smallskip\textit{Formulation and concept references:} \href{https://arxiv.org/pdf/2011.11379}{[S1]}.

\subsection{Short English statement}
Must a compact K\"ahler manifold with no nonconstant entire curves have an ample canonical bundle?
\subsection{Sources}
\begin{itemize}
\item[S1]S. Diverio, Kobayashi hyperbolicity, negativity of the curvature and positivity of the canonical bundle, version 2 (2020).
\url{https://arxiv.org/pdf/2011.11379}.
\textit{Formulation source.}
\end{itemize}
\end{document}
